\scanpage{104}
\mspage{97}

le ``il suffit'' dans b), qui \deleted{\unsure{prouve}} \inserted{\unsure{conclura}} que $X=\varinjlim X_i\in M$, qed.

\uline{NB} On a prouvé $b\Rightarrow a$, sans utiliser l'hyp. \unsure{mise} que tout objet de $\mathrm{Filt}^{\pi}$ est $\pi$-accessible. Mais on a utilisé cette hypothèse pour prouver que b \deleted{implique} (ou a) implique que $M$ est stable par $\varinjlim$ grandes devant $\pi$ (lemme \unsure{2|?}, p. \unsure{90|93}).

$a\overset{?}{\Longrightarrow}b$. \deleted{\unsure{Si}} Le ``il suffit'' de b) est conséquence immédiate de a). Il reste à prouver que si, dans le \deleted{\unsure{théorème} \unsure{de b)}, \unsure{pour}} \inserted{système inductif $\pi$-adapté} $(X_i)_{i\in I}$ de limite $X$, on a $X\in M$, alors $I_M$ est cofinal dans $I$. Mais par hyp. a), on aura $\deleted{\unsure{M}} X=\varinjlim_J Y_j$, avec $J$ grand devant $\pi$, les $Y_j$ dans $M^\pi$. \deleted{Soit $i\in I$,} Admettons que $\mathrm{Filt}^\pi\mathrel{\unsure{\subset|\in}} E_\pi$ i.e tt élément de $\mathrm{Filt}^{\unsure{\pi|n}}$, en particulier les $X_i$, sont \deleted{bien} $\pi$-accessibles.

Je pense que je n'arriverai pas à prouver $a_\pi\Rightarrow b_\pi$, mais seulement ceci : Si $\pi'$ est un cardinal $>\pi$ (disons le successeur de $\pi$ \unsure{$\aleph_{\unsure{\gamma|\eta}}$} dans la \deleted{\unsure{suite}} \inserted{suite ordinale} des cardinaux) alors

\[
 a_\pi\Longrightarrow b_{\pi'}.
\]

\deleted{Donc dans l'argument qui précède, on supposera $I$ grand devant $\pi'$, non seulement grand devant $\pi$. Soit \inserted{alors} $\alpha_0$ le plus petit ordinal qui soit grand devant $\pi$, il est de cardinal $\leq\pi'$ (car comme $\pi'>\pi$, le plus petit ordinal $\alpha'_0$ de cardinal $\pi'$ est grand dev\textsuperscript{ant} $\pi$ car $\pi<\pi'$, donc $\alpha'\geq\alpha'_0$, donc $\pi'=\operatorname{card}\alpha'\geq\operatorname{card}\unsure{\alpha_0|\pi_0}$.}

% archive sheet number: 103
